\documentclass[11pt,a4paper]{article}
\usepackage[italian]{babel}
\usepackage[margin=2.5cm]{geometry}
\usepackage{parskip}
\usepackage{amsmath, amssymb, amsthm}
\usepackage{booktabs}
\usepackage{array}
\usepackage{xcolor}
\usepackage{needspace}
\usepackage{enumitem}
\usepackage{listings}
\usepackage{tikz}
\usepackage[most]{tcolorbox}
\usepackage[hidelinks]{hyperref}
\usetikzlibrary{decorations.pathreplacing, shapes.geometric, arrows.meta, positioning}

% Niente vedove/orfane: i paragrafi non lasciano righe isolate a fine/inizio pagina
\widowpenalty=10000
\clubpenalty=10000

% Elenchi più compatti
\setlist{itemsep=2pt, parsep=0pt, topsep=4pt, partopsep=0pt}

% --- Riquadri con bordo nero, senza numero (si spezzano tra due pagine) ---
% Uso: \begin{definizione} ... \end{definizione}
%      \begin{teorema}[Nome] ... \end{teorema}  (il nome è facoltativo)
\tcbset{nota/.style={enhanced jigsaw, breakable, colback=white,
  colframe=black, boxrule=0.5pt, sharp corners}}
\NewTColorBox{definizione}{}{nota,
  before upper={\textbf{Definizione.}\ }}
\NewTColorBox{teorema}{o}{nota,
  before upper={\textbf{Teorema\IfValueT{#1}{ (#1)}.}\ }}
\NewTColorBox{proposizione}{o}{nota,
  before upper={\textbf{Proposizione\IfValueT{#1}{ (#1)}.}\ }}
\NewTColorBox{proprieta}{o}{nota,
  before upper={\textbf{Proprietà\IfValueT{#1}{ (#1)}.}\ }}
\NewTColorBox{assioma}{o}{nota, boxrule=1pt,
  before upper={\textbf{Assioma\IfValueT{#1}{ (#1)}.}\ }}

% --- Ambienti semplici (senza riquadro) ---
% Il titolo sta in cima al blocco, anche se il contenuto inizia con un riquadro o
% con codice; e non resta mai da solo in fondo a una pagina.
% Uso: \begin{esempio}[Titolo facoltativo] ... \end{esempio}
\newcommand{\newrunin}[2]{%
  \NewDocumentEnvironment{#1}{o}{%
    \par\Needspace{4\baselineskip}\addvspace{\medskipamount}\noindent
    \textbf{#2}\IfValueT{##1}{ (##1)}\textbf{.}\ \ignorespaces}%
  {\par\addvspace{\medskipamount}}}
\newrunin{esempio}{Esempio}
\newrunin{esercizio}{Esercizio}
\newrunin{osservazione}{Osservazione}

% V e F nelle tavole di verità
\newcommand{\V}{\textsf{V}}
\newcommand{\F}{\textsf{F}}

% --- Codice C ---
% Uso: \begin{codice} ... \end{codice}      codice C numerato
%      \begin{comando} ... \end{comando}    comandi da terminale
%      \begin{risultato} ... \end{risultato} output di un programma
%      \lstinline|printf("ciao");|           codice dentro al testo
\lstdefinestyle{c}{
  language=C,
  basicstyle=\ttfamily\small,
  keywordstyle=\color{blue}\bfseries,
  commentstyle=\color{gray}\itshape,
  stringstyle=\color{red!70!black},
  numbers=left,
  numberstyle=\tiny\color{gray},
  numbersep=8pt,
  columns=flexible,
  keepspaces=true,
  breaklines=true,
  showstringspaces=false,
  tabsize=4,
  morekeywords={include, define},
  % lettere accentate nei commenti
  literate={à}{{\`a}}1 {è}{{\`e}}1 {é}{{\'e}}1 {ì}{{\`i}}1
           {ò}{{\`o}}1 {ù}{{\`u}}1 {È}{{\`E}}1 {À}{{\`A}}1
}
\lstset{style=c}
\newtcblisting{codice}{listing only, nota, left=6mm,
  listing options={style=c}}
\newtcblisting{comando}{listing only, nota,
  listing options={basicstyle=\ttfamily\small, numbers=none, language={},
    columns=flexible, keepspaces=true, breaklines=true}}
\newtcblisting{risultato}{listing only, enhanced jigsaw, breakable,
  colback=black!5, colframe=black, boxrule=0.5pt, sharp corners,
  listing options={basicstyle=\ttfamily\small, numbers=none, language={},
    columns=flexible, keepspaces=true, breaklines=true}}

% --- Diagrammi di flusso (simboli standard) ---
\tikzset{
  terminale/.style = {ellipse, draw, minimum width=2.5cm, minimum height=0.9cm},
  io/.style        = {trapezium, trapezium left angle=70, trapezium right angle=110,
                      draw, minimum width=2.5cm, minimum height=0.9cm},
  processo/.style  = {rectangle, draw, minimum width=2.5cm, minimum height=0.9cm},
  decisione/.style = {diamond, draw, aspect=2, minimum width=2.5cm},
  freccia/.style   = {thick, -{Stealth}}
}

\title{Sistemi di riferimento}
\author{Marco Cerbella}
\date{Lezione 3 -- 1 ottobre 2026}

\begin{document}
\maketitle

\section{Sistemi di riferimento}

Stabilire un sistema di riferimento in uno spazio significa fissare
\begin{enumerate}
    \item un punto detto \emph{origine};
    \item un metodo per determinare univocamente la posizione di tutti gli altri punti rispetto all'origine, associando loro dei numeri reali detti \emph{coordinate}.
\end{enumerate}
La \emph{dimensione} dello spazio è il numero di coordinate necessarie a descrivere la posizione di un punto. Si considerano sistemi unidimensionali (1D), bidimensionali (2D) e tridimensionali (3D).

\section{Sistema unidimensionale}

Un sistema di riferimento unidimensionale è costituito da:
\begin{itemize}
    \item una retta orizzontale detta \emph{asse delle ascisse};
    \item un'origine, di solito indicata con $O$;
    \item un verso di percorrenza, solitamente verso destra;
    \item un punto $U$, distinto da $O$ e situato alla destra di $O$, che fissa come unità di misura la lunghezza del segmento $OU$, indicata con $|OU|$.
\end{itemize}
L'\emph{ascissa} è il numero reale associato ad ogni punto $P$; $|x|$ è la lunghezza del segmento $OP$ secondo l'unità di misura stabilita. Se $x > 0$, $P$ è a destra di $O$; se $x < 0$ è a sinistra; se $x = 0$, $P$ coincide con $O$.

\begin{center}
\begin{tikzpicture}
  \draw[->] (-3,0) -- (4,0);
  \fill (0,0) circle (1.5pt) node[below] {$O$};
  \fill (1,0) circle (1.5pt) node[below] {$U$};
  \fill (2.7,0) circle (1.5pt) node[below] {$P$};
  \node[above] at (2.7,0.05) {$x_P$};
\end{tikzpicture}
\end{center}

\section{Sistema bidimensionale: il piano cartesiano}

Sistema di coordinate cartesiane bidimensionale (piano cartesiano monometrico):
\begin{enumerate}
    \item si fissano nel piano due rette perpendicolari che si intersecano in un punto $O$;
    \item si fissano sulla prima un punto $U$ e sulla seconda un punto $V$, entrambi distinti da $O$, con $|OU| = |OV|$ (unità di misura);
    \item dato un punto $P$ si indicano con $P_1$ e $P_2$ le sue proiezioni ortogonali sulle due rette;
    \item ognuna delle due rette è un sistema di coordinate unidimensionale: si indicano con $x_P$ e $y_P$ le coordinate di $P_1$ e $P_2$.
\end{enumerate}
Ad ogni punto $P$ del piano è così associata una coppia ordinata $(x_P, y_P)$ di numeri reali, le \emph{coordinate cartesiane bidimensionali} di $P$.

Rappresentando il piano, la prima retta è orizzontale (\emph{asse delle $x$} o delle ascisse) e la seconda verticale (\emph{asse delle $y$} o delle ordinate). Per $P = (x_P, y_P)$: $x_P$ è l'ascissa e $y_P$ è l'ordinata di $P$.

Il piano risulta suddiviso in quattro \emph{quadranti} numerati in senso antiorario a partire da quello in cui entrambe le coordinate sono positive.

\begin{center}
\begin{tikzpicture}[scale=0.9]
  \draw[->] (-3.2,0) -- (3.5,0) node[right] {$x$};
  \draw[->] (0,-2.5) -- (0,3) node[above] {$y$};
  \fill (2,1.5) circle (1.5pt) node[above right] {$P = (x_P, y_P)$};
  \draw[dashed] (2,0) -- (2,1.5);
  \draw[dashed] (0,1.5) -- (2,1.5);
  \node[below] at (2,0) {$x_P$};
  \node[left] at (0,1.5) {$y_P$};
  \node at (1.3,2.4) {I}; \node at (-1.3,2.4) {II};
  \node at (-1.3,-1.6) {III}; \node at (1.3,-1.6) {IV};
\end{tikzpicture}
\end{center}

\section{Sistema tridimensionale: lo spazio cartesiano}

Sistema di coordinate cartesiane tridimensionale (spazio cartesiano monometrico):
\begin{enumerate}
    \item si fissano nello spazio tre rette perpendicolari che si intersecano in un punto $O$ detto origine;
    \item si fissano sulla prima un punto $U$, sulla seconda un punto $V$ e sulla terza un punto $W$, tutti distinti da $O$, con $|OU| = |OV| = |OW|$;
    \item dato un punto $P$ si indicano con $P_1, P_2, P_3$ le proiezioni ortogonali sulle tre rette;
    \item ognuna delle rette è un sistema di coordinate unidimensionale: si indicano con $x_P, y_P, z_P$ le coordinate di $P_1, P_2, P_3$.
\end{enumerate}
Ad ogni punto dello spazio è associata una terna ordinata $(x_P, y_P, z_P)$ di numeri reali, le \emph{coordinate cartesiane tridimensionali} di $P$:
\begin{itemize}
    \item \emph{ascissa}: primo numero della terna, sull'asse $x$;
    \item \emph{ordinata}: secondo numero della terna, sull'asse $y$;
    \item \emph{quota}: terzo numero della terna, sull'asse $z$.
\end{itemize}

\begin{center}
\begin{tikzpicture}[scale=1]
  \draw[->] (0,0) -- (0,3) node[above] {$z$};
  \draw[->] (0,0) -- (3.2,0) node[right] {$y$};
  \draw[->] (0,0) -- (-1.6,-1.2) node[below left] {$x$};
  \fill (1.6,1.5) circle (1.5pt) node[above right] {$P = (x_P, y_P, z_P)$};
\end{tikzpicture}
\end{center}

\end{document}
